\section{Demostración de Pertenencia a NP}

Para demostrar que la versión de decisión del \textit{Hitting-Set Problem} pertenece a la clase NP, debemos probar que dada una solución candidata (un certificado), es posible verificar su validez en un tiempo acotado por un polinomio.

A continuación, se presenta el algoritmo certificador escrito en Python que valida una solución propuesta:

\begin{lstlisting}[language=Python, frame=single, numbers=left, breaklines=true, basicstyle=\ttfamily\footnotesize]
def verificador_hitting_set(k, solucion, subsets, universo):
    """
    Verifica en tiempo polinomial si una solución candidata 
    es válida para el Hitting Set Problem (versión decisión).
    """
    
    # 1. Verificar restricción de cardinalidad
    if len(solucion) > k:
        return False
        
    # 2. Verificar que los elementos pertenezcan al universo
    for elemento in solucion:
        if elemento not in universo:
            return False
            
    # 3. Verificar condición de intersección (Hitting Set)
    for subconjunto in subsets:
        interseccion_encontrada = False
        
        for elemento in solucion:
            if elemento in subconjunto:
                interseccion_encontrada = True
                break # Globo pinchado, pasamos al siguiente
                
        if not interseccion_encontrada:
            return False
            
    return True
\end{lstlisting}

\subsection*{Análisis de Complejidad}

Para justificar que el validador opera en tiempo polinomial, definimos las variables que representan el tamaño de nuestra entrada:
\begin{itemize}
    \item $n$: Cantidad de elementos en el \texttt{universo} ($|A|$).
    \item $m$: Cantidad de subconjuntos en \texttt{subsets}.
    \item $c$: Cantidad de elementos en la \texttt{solucion} candidata ($|C|$). Por definición lógica del problema, $c \leq n$.
\end{itemize}

Evaluamos el costo temporal de cada paso del algoritmo en el peor de los casos:
\begin{enumerate}
    \item \textbf{Verificación de cardinalidad (Línea 8):} Obtener la longitud de la solución candidata y compararla con $k$ toma un tiempo de $\mathcal{O}(1)$ (o a lo sumo $\mathcal{O}(c)$ dependiendo de la implementación subyacente de la estructura de datos).
    \item \textbf{Pertenencia al universo (Líneas 12-14):} Se itera sobre los $c$ elementos de la solución, verificando para cada uno si pertenece al universo de tamaño $n$. El costo en el peor caso es $\mathcal{O}(c \cdot n)$.
    \item \textbf{Condición de Hitting Set (Líneas 17-26):} Se itera sobre los $m$ subconjuntos. Para cada uno (que puede tener hasta $n$ elementos), se busca si contiene alguno de los $c$ elementos de la solución. Este proceso de búsquedas anidadas requiere a lo sumo $\mathcal{O}(m \cdot c \cdot n)$ operaciones.
\end{enumerate}

El tiempo total de ejecución en el peor de los casos está acotado asintóticamente por el paso más costoso: $\mathcal{O}(m \cdot c \cdot n)$. Dado que la cantidad de elementos en la solución nunca será mayor al universo total ($c \leq n$), podemos simplificar la cota superior a $\mathcal{O}(m \cdot n^2)$.

Como la función $\mathcal{O}(m \cdot n^2)$ es un polinomio respecto al tamaño de la entrada, queda formalmente demostrado que la verificación de un certificado se puede realizar en tiempo polinomial. En conclusión, \textbf{el Hitting-Set Problem se encuentra en la clase NP}.